\documentclass[11pt]{article}
\usepackage[margin=1in]{geometry}
\usepackage{booktabs}
\usepackage{array}
\usepackage{xcolor}
\usepackage{amsmath}
\usepackage{hyperref}
\usepackage{longtable}
\usepackage{graphicx}

\newcommand{\ratebar}[1]{\textcolor{blue!65!black}{\rule{#1mm}{6pt}}}

\hypersetup{
  colorlinks=true,
  linkcolor=blue,
  urlcolor=blue
}

\title{Archive-2 Dataset Audit and Clean Ingestion Report}
\author{TP2-Inspired SPY Call Option Project}
\date{June 8, 2026}

\begin{document}
\maketitle

\section*{Purpose}

This report documents the first integration step for the professor-provided
\texttt{Archive-2} SPY option dataset. The goal was to inspect whether the new
2013--2026 dataset is suitable for the publishable call-only American-option
research project, identify real data caveats, and create a clean derived dataset
without modifying the raw professor-provided files.

\section*{June 10 Restart: Active Clean Data Layer}

After the raw EDA, we changed the active clean-data assumption. The earlier
derived datasets required \texttt{quote\_time >= 15:55 ET}. That rule was too
restrictive for this archive because the API appears to store the latest
available scraped quote for many contracts even when the timestamp is earlier
than 15:55. Therefore, the next pipeline rebuild will use the latest available
quote per contract-day after duplicate collapse, not a hard 15:55 cutoff.

The new active clean file is:

\begin{quote}
\small\path{data/processed/archive2_spy_call_allmoneyness_lastquote_clean_2013_2026.parquet}
\end{quote}

This file is derived from the raw professor-provided
\texttt{Archive-2/spy\_options\_daily/year=*/data.parquet} files. The raw files
remain unmodified. The cleaning rules are:

\begin{enumerate}
  \item Collapse raw duplicates to one row per \texttt{quote\_date,
  contract\_id}, keeping the latest available quote fields and summing duplicate
  volume rows.
  \item Keep SPY calls only across all moneyness levels.
  \item Keep DTE from 1 to 90 calendar days.
  \item Require bid $>$ 0, ask $>$ bid, and mid $>$ 0.
  \item Require spread percentage $\leq 50\%$.
  \item Do not require \texttt{quote\_time >= 15:55}; the retained row is the
  latest available quote for that contract-day.
\end{enumerate}

\begin{center}
\begin{tabular}{lrrrrrr}
\toprule
Dataset Version & Rows & Sessions & Duplicate Rows & Call Rows & Min DTE & Max DTE \\
\midrule
Old near-close $\geq$15:55 & 804,809 & 3,255 & 0 & 804,809 & 1 & 90 \\
New last-available quote & 4,912,891 & 3,256 & 0 & 4,912,891 & 1 & 90 \\
\bottomrule
\end{tabular}
\end{center}

The session count increases by only one date because most dates already had at
least one valid 15:55-or-later quote. The row count increases substantially
because many individual option contracts have their latest available quote
before 15:55 or have a missing quote timestamp but valid bid/ask prices. The
one added session is November 25, 2022, an early-close style date whose retained
quotes end around 13:15 ET.

\subsection*{Pipeline To Rebuild From This Data}

The previous day-by-day pipeline is now treated as a prototype. The next build
will reconstruct the research path from the active last-available dataset in
the following order:

\begin{enumerate}
  \item Run forward-adjusted TP2 detection on the new all-call clean panel.
  \item Create a stable date-independent rectangle key from strikes,
  expirations, and A/B/C/D option symbols.
  \item Track violation episodes across consecutive trading dates so the same
  continuing rectangle is counted once, not once per day.
  \item Label episode starts for reversion, time-to-reversion, fixed-horizon
  P\&L, and flexible-exit P\&L.
  \item Expand episode starts into the five strategy families:
  \texttt{full4}, \texttt{T1}, \texttt{T2}, \texttt{K1}, and \texttt{K2}.
  \item Train ML only on point-in-time episode-start features using a
  leakage-safe rolling walk-forward split by episode start date.
  \item Backtest one entry per selected episode with invalid-selected accounting,
  integer sizing, scale-factor caps, liquidity controls, and transaction costs.
\end{enumerate}

The main methodological correction is episode tracking. If the same TP2
rectangle violates on multiple consecutive trading days, it should be one
economic violation episode until it reverts, expires, disappears, or becomes
untradeable. This prevents inflated violation counts and repeated entries into
the same continuing mispricing.

\subsection*{First Smoke Test on New Active Data}

A four-session forward-adjusted TP2 detection smoke test was run on the new
last-available all-call dataset for January 2--5, 2024:

\begin{quote}
\small\path{data/features/archive2_lastquote_smoke/tp2_call_forward_features_2024_0102_0105.parquet}
\end{quote}

\begin{center}
\begin{tabular}{rrrrr}
\toprule
Input Rows & Sessions & Violation Rows & Sessions with Violations & Duplicate IDs \\
\midrule
7,232 & 4 & 661 & 4 & 0 \\
\bottomrule
\end{tabular}
\end{center}

This confirms that the existing forward-adjusted TP2 detector can consume the
new last-available quote panel. The smoke output spans four trading sessions,
has 122 feature columns, and has normalized severity ranging from 0.000103 to
0.285714. No ML, label construction, or backtest has been rerun yet on the new
active data.

\section*{Current Pipeline Architecture}

\begin{center}
\includegraphics[width=0.98\textwidth]{figures/archive2_pipeline_architecture_detailed.png}
\end{center}

\section*{Raw Dataset Summary}

The new archive is substantially stronger than the previous JSON archive. It
contains SPY option data from April 1, 2013 through June 4, 2026, with yearly
Parquet partitions. The documented source is Databento OPRA consolidated NBBO
data, with raw SPY close, dividend fields, computed forward/rate fields, and
Black--Scholes-derived IV/Greeks.

\begin{longtable}{>{\raggedright\arraybackslash}p{0.25\textwidth}p{0.65\textwidth}}
\toprule
\textbf{Field Group} & \textbf{Available in Archive-2} \\
\midrule
Option identifiers & \texttt{contract\_id}, \texttt{option\_symbol}, expiration, strike, call/put flag \\
Prices & bid, ask, mid \\
Tradability & bid size, ask size, daily volume, open interest \\
Underlying & raw \texttt{spy\_close} \\
Dividend controls & next ex-dividend date and dividend amount \\
Model features & implied volatility, delta, theta, vega \\
Forward/rate inputs & parity-implied forward and risk-free rate \\
\bottomrule
\end{longtable}

\section*{Caveats Found}

Three caveats were audited before ingestion.

\subsection*{Skipped Sessions}

\texttt{skipped\_days.txt} contains 56 skipped sessions. Of these, 28 are
known early-close or half-day style dates such as July 3, Black Friday, and
Christmas Eve. The remaining 28 dates are normal weekdays in 2024--2025 and
should be treated as recoverable data gaps, not market-calendar effects. The
permanent vendor gap documented in \texttt{known\_gaps.txt} is June 18, 2021.

\subsection*{Quote Freshness}

The quote-time caveat is real. Normal full trading days contain many option rows
whose \texttt{quote\_time} is much earlier than the market close, including
morning and midday timestamps. This is not only a half-day effect. For a
publishable TP2-style study, all legs in a detected rectangle should use fresh
near-close quotes. The derived dataset therefore applies a closing freshness
filter:

\[
\texttt{quote\_time} \geq \text{15:55 ET}.
\]

A stricter 15:58 ET filter can be used later as a robustness check.

\subsection*{Duplicate Rows}

The raw yearly files contain repeated rows for the same
\texttt{quote\_date, contract\_id}. This appears to come from merging
unaggregated daily volume records into quote rows. The quote fields are often
identical, while the volume field can differ. This is not a natural market
caveat; it is a data-construction issue. The clean ingestion step collapses to
one row per \texttt{quote\_date, contract\_id} and sums duplicate volume rows.

\section*{Clean Ingestion Added}

Two new project files were added:

\begin{itemize}
  \item \texttt{src/tp2rr2/archive2\_ingest.py}
  \item \texttt{scripts/run\_ingest\_archive2.py}
\end{itemize}

The ingestion script treats \texttt{Archive-2} as immutable raw data. It writes
cleaned outputs under \texttt{data/processed/}; no professor-provided raw files
are overwritten or modified.

\section*{Cleaning Rules}

The derived research dataset applies the following filters:

\begin{enumerate}
  \item Deduplicate to one row per \texttt{quote\_date + contract\_id}.
  \item Sum duplicate daily volume rows.
  \item Normalize column names for the existing project pipeline.
  \item Keep calls only.
  \item Keep out-of-the-money calls only: strike $>$ raw SPY close.
  \item Keep DTE from 1 to 90 calendar days.
  \item Require bid $>$ 0, ask $>$ bid, and mid $>$ 0.
  \item Require spread percentage $\leq 50\%$.
  \item Require fresh quote time at or after 15:55 ET.
\end{enumerate}

\section*{Cleaned Output}

The main cleaned file is:

\begin{center}
\texttt{data/processed/archive2\_spy\_call\_otm\_clean\_2024\_2026.parquet}
\end{center}

The summary file is:

\begin{center}
\texttt{data/processed/archive2\_spy\_call\_otm\_clean\_2024\_2026\_summary.json}
\end{center}

A full-history cleaned file was later created using the same cleaning rules:

\begin{center}
\texttt{data/processed/archive2\_spy\_call\_otm\_clean\_2013\_2026.parquet}
\end{center}

with summary:

\begin{center}
\texttt{data/processed/archive2\_spy\_call\_otm\_clean\_2013\_2026\_summary.json}
\end{center}

\begin{table}[h]
\centering
\begin{tabular}{rrrrr}
\toprule
Year & Raw Rows & Deduped Rows & Clean Rows & Clean Sessions \\
\midrule
2024 & 6,044,559 & 2,147,976 & 55,799 & 236 \\
2025 & 7,119,099 & 2,256,872 & 65,726 & 232 \\
2026 & 3,677,988 & 1,314,233 & 29,562 & 103 \\
\midrule
Total & 16,841,646 & 5,719,081 & 151,087 & 571 \\
\bottomrule
\end{tabular}
\caption{Raw, deduplicated, and cleaned row counts for the 2024--2026 call-only OTM research sample.}
\end{table}

For the full 2013--2026 sample, the same ingestion rules produced 536,611
cleaned rows across 3,255 sessions from 58,834,075 raw rows and 22,589,850
deduplicated contract-date rows. The full-history cleaned dataset passed the
same validation checks: calls only, OTM only, DTE 1--90, no duplicate
\texttt{date, contractID} rows, valid bid/ask, and quote times from
15:55:00--15:59:59 ET.

\section*{Archive-2 Pipeline Files}

The current Archive-2 project pipeline is separated from the older SPY JSON
pipeline. The files below are part of the new call-only research path:

\begin{longtable}{>{\raggedright\arraybackslash}p{0.35\textwidth}p{0.55\textwidth}}
\toprule
\textbf{File} & \textbf{Role} \\
\midrule
\texttt{src/tp2rr2/archive2\_ingest.py} & Clean ingestion utilities for the professor-provided Archive-2 panel. \\
\texttt{scripts/run\_ingest\_archive2.py} & Command-line runner that creates the cleaned OTM call dataset. \\
\texttt{data/processed/archive2\_spy\_call\_otm\_clean\_2024\_2026.parquet} & Cleaned call-only OTM dataset used by the new pipeline. \\
\texttt{data/processed/archive2\_spy\_call\_otm\_clean\_2024\_2026\_summary.json} & Cleaning summary and row-count audit. \\
\texttt{data/processed/archive2\_spy\_call\_otm\_clean\_2013\_2026.parquet} & Full-history cleaned call-only OTM dataset using the same rules. \\
\texttt{data/processed/archive2\_spy\_call\_otm\_clean\_2013\_2026\_summary.json} & Full-history cleaning summary and row-count audit. \\
\texttt{src/tp2rr2/detect\_forward\_tp2.py} & Forward-adjusted TP2 call detector using rounded adjusted strikes. \\
\texttt{scripts/run\_detect\_archive2\_forward.py} & Command-line runner for forward-adjusted TP2 detection and feature generation. \\
\texttt{data/features/archive2\_forward/tp2\_call\_forward\_features\_2024\_2026.parquet} & Forward-adjusted TP2 call violation feature table. \\
\texttt{data/features/archive2\_forward/tp2\_call\_forward\_summary\_2024\_2026.json} & Forward-adjusted detection summary. \\
\texttt{data/features/archive2\_forward/tp2\_call\_forward\_features\_2013\_2026.parquet} & Full-history forward-adjusted TP2 call violation feature table. \\
\texttt{data/features/archive2\_forward/tp2\_call\_forward\_summary\_2013\_2026.json} & Full-history forward-adjusted detection summary. \\
\texttt{reports/archive2\_dataset\_ingestion\_report.tex} & Source for this report. \\
\texttt{reports/archive2\_dataset\_ingestion\_report.pdf} & PDF version of this report. \\
\bottomrule
\end{longtable}

\section*{Forward-Adjusted TP2 Detection}

A main-specification forward-adjusted TP2 detector was added after the clean
ingestion step. It implements the call inequality motivated by
Glasserman--Li--Pirjol:

\[
C(K_1,T_1)C(K_2,T_2)
\geq
C(K_1F_{T_2}/F_{T_1},T_2)C(K_2F_{T_1}/F_{T_2},T_1),
\]

for $T_1 < T_2$ and $K_1/F_{T_1} < K_2/F_{T_2}$. The adjusted strikes are
rounded up to the nearest available listed strike in the relevant expiration.
This keeps the two-strike approximation available as a robustness design, while
making the forward-adjusted condition the primary detection method for the
publishable Archive-2 pipeline.

\begin{table}[h]
\centering
\begin{tabular}{rrrr}
\toprule
Year & Violation Rows & Sessions with Violations & Median Normalized Severity \\
\midrule
2024 & 13,643 & 182 & 0.087329 \\
2025 & 37,910 & 194 & 0.097852 \\
2026 & 9,054 & 88 & 0.111111 \\
\midrule
Total & 60,607 & 464 & 0.097473 \\
\bottomrule
\end{tabular}
\caption{Forward-adjusted TP2 call violations detected on the cleaned 2024--2026 Archive-2 sample.}
\end{table}

The same detector was then run on the full 2013--2026 cleaned dataset.

\begin{table}[h]
\centering
\begin{tabular}{rrrr}
\toprule
Period & Violation Rows & Sessions with Violations & Median Normalized Severity \\
\midrule
2013--2026 & 109,160 & 1,286 & 0.075509 \\
2024--2026 & 60,607 & 464 & 0.097473 \\
\bottomrule
\end{tabular}
\caption{Forward-adjusted TP2 call detection totals for full-history and recent samples.}
\end{table}

The full-history output has no duplicate \texttt{violation\_id} values and spans
June 20, 2013 through June 4, 2026. Violations are sparse in early years and
increase substantially in 2020 and again in 2024--2026, which should be treated
as an empirical finding to analyze later rather than assumed to be a model edge.

\subsection*{Yearly Violation Percentages}

The Glasserman--Li--Pirjol-style violation rate is the fraction of candidate
strike-expiration rectangles that violate TP2, not the fraction of sessions that
contain at least one violation. This distinction is important because each
session can contain thousands of candidate rectangles; a low per-candidate
violation rate can still imply that many sessions contain at least one
violation.

The table below reports the corrected candidate-rectangle violation rate by
year:

\begin{center}
\begin{tabular}{rrrrl}
\toprule
Year & Candidate Rectangles & Violations & Candidate Violation \% & Visual Bar \\
\midrule
2013 & 38,899 & 19 & 0.0488\% & \ratebar{1.9} \\
2014 & 75,401 & 30 & 0.0398\% & \ratebar{1.5} \\
2015 & 198,179 & 146 & 0.0737\% & \ratebar{2.8} \\
2016 & 224,309 & 282 & 0.1257\% & \ratebar{4.8} \\
2017 & 102,245 & 130 & 0.1271\% & \ratebar{4.8} \\
2018 & 658,754 & 1,510 & 0.2292\% & \ratebar{8.7} \\
2019 & 613,731 & 292 & 0.0476\% & \ratebar{1.8} \\
2020 & 3,531,543 & 26,653 & 0.7547\% & \ratebar{28.7} \\
2021 & 1,787,279 & 3,526 & 0.1973\% & \ratebar{7.5} \\
2022 & 4,946,796 & 10,905 & 0.2204\% & \ratebar{8.4} \\
2023 & 2,313,380 & 5,060 & 0.2187\% & \ratebar{8.3} \\
2024 & 2,338,601 & 13,643 & 0.5834\% & \ratebar{22.2} \\
2025 & 3,568,863 & 37,910 & 1.0622\% & \ratebar{40.4} \\
2026 & 1,502,185 & 9,054 & 0.6027\% & \ratebar{22.9} \\
\bottomrule
\end{tabular}
\end{center}

These corrected rates are consistent with the intuition that TP2 violations are
rare at the candidate-pair level. The earlier session-incidence view is still
useful operationally, because it tells us how often the strategy sees at least
one possible signal, but it should not be called the theoretical violation rate.

For reference, the next table reports the percentage of cleaned trading sessions
in each year with at least one forward-adjusted TP2 call violation.

\begin{center}
\begin{tabular}{rrrrl}
\toprule
Year & Clean Sessions & Violation Sessions & Session \% & Visual Bar \\
\midrule
2013 & 189 & 8 & 4.23\% & \ratebar{2.1} \\
2014 & 249 & 12 & 4.82\% & \ratebar{2.4} \\
2015 & 250 & 31 & 12.40\% & \ratebar{6.2} \\
2016 & 251 & 48 & 19.12\% & \ratebar{9.6} \\
2017 & 249 & 21 & 8.43\% & \ratebar{4.2} \\
2018 & 248 & 58 & 23.39\% & \ratebar{11.7} \\
2019 & 249 & 33 & 13.25\% & \ratebar{6.6} \\
2020 & 251 & 123 & 49.00\% & \ratebar{24.5} \\
2021 & 250 & 150 & 60.00\% & \ratebar{30.0} \\
2022 & 250 & 184 & 73.60\% & \ratebar{36.8} \\
2023 & 248 & 154 & 62.10\% & \ratebar{31.1} \\
2024 & 236 & 182 & 77.12\% & \ratebar{38.6} \\
2025 & 232 & 194 & 83.62\% & \ratebar{41.8} \\
2026 & 103 & 88 & 85.44\% & \ratebar{42.7} \\
\bottomrule
\end{tabular}
\end{center}

\begin{center}
\begin{tabular}{rrr}
\toprule
Year & Violation Rows & Violation Rows / Clean Rows \\
\midrule
2013 & 19 & 0.24\% \\
2014 & 30 & 0.27\% \\
2015 & 146 & 0.84\% \\
2016 & 282 & 1.46\% \\
2017 & 130 & 0.90\% \\
2018 & 1,510 & 5.03\% \\
2019 & 292 & 0.93\% \\
2020 & 26,653 & 41.57\% \\
2021 & 3,526 & 7.00\% \\
2022 & 10,905 & 13.51\% \\
2023 & 5,060 & 8.63\% \\
2024 & 13,643 & 24.45\% \\
2025 & 37,910 & 57.68\% \\
2026 & 9,054 & 30.63\% \\
\bottomrule
\end{tabular}
\end{center}

The source tables for these rates are saved at
\texttt{reports/figures/archive2\_forward\_candidate\_violation\_rates\_by\_year.csv}
and
\texttt{reports/figures/archive2\_forward\_violation\_rates\_by\_year.csv}.

\subsection*{Yearly Chain Breadth}

The number of TP2 candidate rectangles depends strongly on option-chain breadth:
how many strikes and expirations are available each session. The cleaned
OTM-call detection universe shows that recent years contain substantially more
listed structure than the early sample. This helps explain why 2024--2025 can
produce more selected TP2 trades even when many individual contracts have low
or zero reported daily volume.

\begin{center}
\begin{tabular}{rrrrrr}
\toprule
Year & Sessions & Unique Strikes & Unique Expiries & Avg Strikes/Day & Avg Contracts/Day \\
\midrule
2013 & 189 & 79 & 49 & 14.1 & 41.3 \\
2014 & 249 & 94 & 66 & 14.5 & 45.4 \\
2015 & 250 & 83 & 63 & 21.8 & 69.7 \\
2016 & 251 & 121 & 84 & 22.5 & 76.7 \\
2017 & 249 & 108 & 118 & 17.3 & 58.2 \\
2018 & 248 & 143 & 169 & 28.0 & 121.1 \\
2019 & 249 & 191 & 174 & 25.4 & 126.0 \\
2020 & 251 & 269 & 178 & 43.4 & 255.4 \\
2021 & 250 & 183 & 176 & 35.6 & 201.6 \\
2022 & 250 & 204 & 184 & 54.2 & 323.0 \\
2023 & 248 & 156 & 266 & 38.4 & 236.3 \\
2024 & 236 & 225 & 268 & 36.6 & 236.4 \\
2025 & 232 & 289 & 268 & 46.6 & 283.3 \\
2026 & 103 & 199 & 123 & 48.4 & 287.0 \\
\bottomrule
\end{tabular}
\end{center}

The source table is saved at
\texttt{reports/tables/archive2\_otm\_call\_chain\_breadth\_by\_year.csv}.

Validation checks on the forward-adjusted output:

\begin{itemize}
  \item 60,607 feature rows were produced from 151,087 cleaned option rows.
  \item 464 sessions contain at least one forward-adjusted TP2 call violation.
  \item No duplicate \texttt{violation\_id} values were found.
  \item The output includes the forward prices, raw adjusted strikes, rounded
  adjusted strikes, leg-level quote fields, and existing project feature fields.
\end{itemize}

\section*{Validation Results}

The cleaned file passed the following checks:

\begin{itemize}
  \item 151,087 rows across 571 sessions.
  \item No duplicate \texttt{date, contractID} rows.
  \item Date range: January 2, 2024 through June 4, 2026.
  \item Calls only.
  \item DTE range: 1--90.
  \item All rows are OTM calls.
  \item Bid/ask validity holds for all rows.
  \item Maximum spread percentage is 50\%.
  \item Quote-time range is 15:55:00--15:59:59 ET.
\end{itemize}

\section*{Post-Ingestion Recheck}

A follow-up audit was run after creating the cleaned dataset.

\begin{itemize}
  \item No raw \texttt{Archive-2} columns were lost. The cleaned file preserves all
  original fields and adds normalized project columns such as \texttt{date},
  \texttt{expiration}, \texttt{contractID}, \texttt{type},
  \texttt{underlying\_price}, \texttt{spread}, and \texttt{spread\_pct}.
  \item Important model and control fields survived ingestion: bid, ask, mid,
  bid/ask size, volume, open interest, raw SPY close, dividend fields, forward,
  risk-free rate, implied volatility, delta, theta, and vega.
  \item The cleaned file has no duplicate \texttt{date, contractID} rows.
  \item The current TP2 detector successfully ran on a four-session smoke test
  using the cleaned file and produced 238 call-violation feature rows.
  \item The 15:55 quote-freshness filter is restrictive but not destructive. It
  keeps 151,087 rows across all 571 available sessions in 2024--2026.
\end{itemize}

One caution remains: duplicate raw rows often differ only in the daily-volume
field. The ingestion currently sums duplicate volume values because they appear
to originate from unaggregated daily-volume records. For liquidity filtering,
positive-volume indicators are reliable enough for the next stage; exact volume
magnitudes should be interpreted cautiously until the duplicate-volume source is
fully validated against vendor documentation or a direct Databento query.

Additional full-history notes:

\begin{itemize}
  \item 2023 has much weaker positive-open-interest coverage after filtering
  than surrounding years and should be checked before using OI as a strict model
  gate.
  \item 2026 dividend fields are only partly populated because future realized
  dividend dates and amounts are not fully known in the source fields. For
  point-in-time backtests, dividend features should be handled as risk controls,
  not as predictive look-ahead inputs.
\end{itemize}

\section*{Project Impact}

The new dataset is suitable for the publishable American-style SPY call-option
project after cleaning. The main research framing should remain:

\begin{quote}
We use TP2 theory for European call prices as motivation, then study
TP2-inspired relative-value violations in American-style SPY calls with
dividend, liquidity, quote freshness, and after-cost trading controls.
\end{quote}

The current cleaned sample is appropriate for the next pipeline stage:
forward-adjusted TP2 detection, two-strike robustness, label construction, and
machine-learning-based trade selection.

\section*{June 8 Full Pipeline Recheck}

A full recheck was run from raw data assumptions through the current
forward-adjusted feature table.

\begin{itemize}
  \item The professor-provided raw \texttt{Archive-2} files remain unmodified.
  \item The old JSON archive was removed from the active path; the new pipeline
  uses \texttt{Archive-2} Parquet data.
  \item The full 2013--2026 cleaned file has 536,611 rows across 3,255 sessions,
  no duplicate \texttt{date, contractID} rows, calls only, OTM calls only, DTE
  1--90, valid bid/ask, maximum spread percentage 50\%, and no quotes earlier
  than 15:55 in the retained sample.
  \item The full forward-adjusted feature file has 109,160 violation rows across
  1,286 sessions, no duplicate \texttt{violation\_id} values, valid
  $T_1<T_2$, valid forward-moneyness ordering, and no failed TP2 inequality
  checks among retained violations.
  \item Every A/B/C/D entry leg in the forward-adjusted feature table maps back
  to a retained cleaned quote on the same entry date.
  \item The corrected candidate-rectangle violation percentages remain the
  right theoretical rate to report; session incidence should be presented only
  as an operational signal-frequency statistic.
\end{itemize}

The recheck also found one important next-stage issue. A label-construction
smoke test on 100 recent forward-adjusted violations produced only 4 valid
strategy labels out of 500 expanded strategy rows when using the OTM-only
cleaned file as the exit quote lookup. This is expected but not acceptable for
final research labels: a contract that is OTM at signal time can become ATM or
ITM by the exit date and therefore disappear from an OTM-only exit universe.
The detection universe should remain OTM call focused, but the exit-quote panel
must be broader: fresh, valid SPY call quotes across all moneyness levels, still
deduplicated and bid/ask filtered.

\subsection*{Step 1: Broad Exit Quote Panel}

A broad call-only exit quote panel was created for label construction:

\begin{center}
\texttt{data/processed/archive2\_spy\_call\_allmoneyness\_clean\_2013\_2026.parquet}
\end{center}

This file uses the same raw \texttt{Archive-2} source and the same quote-quality
rules as the OTM detection dataset, but it does not apply the OTM-only filter.
It keeps valid SPY calls across all moneyness levels. The output has 804,809
rows across 3,255 sessions, no duplicate \texttt{date, contractID} rows, calls
only, DTE 1--90, valid bid/ask, maximum spread percentage 50\%, and quote times
at or after 15:55 in the retained sample. It contains the original 536,611 OTM
call rows plus 268,198 ATM/ITM call rows for exit pricing.

\subsection*{Step 2: Stable Option Symbol for Exit Lookup}

The first broad-panel label smoke test still produced too few valid labels.
The cause was not the TP2 detector: \texttt{contract\_id} changed across days for
the same listed option in the professor-provided data. For example, the same
OSI option symbol \texttt{SPY 240108C00471000} had different contract IDs on
January 4 and January 5, 2024. The label pipeline was therefore patched to carry
\texttt{option\_symbol} on each A/B/C/D leg and to use the option symbol as the
primary exit lookup key, with contract ID retained only as a fallback.

After this fix, the 100-violation smoke test improved from 10 valid strategy
rows to 215 valid strategy rows out of 500. The remaining invalid labels are
mostly genuine target-date quote availability or short-expiration cases.

\subsection*{Step 3: Full 2013--2026 Labels}

The full forward-adjusted feature table was regenerated with A/B/C/D option
symbols. The violation count remained unchanged at 109,160 rows, confirming
that the metadata fix did not alter detection math.

Full strategy labels were then built using OTM-entry features and the broad
all-moneyness call exit panel:

\begin{center}
\texttt{data/labels/archive2\_forward\_strategy\_training\_rows\_2013\_2026.parquet}
\end{center}

The label file has 545,800 strategy rows, equal to 109,160 violations times
five strategy variants. Of these, 105,231 rows have valid realized-return labels
under the current 5-trading-day/pre-expiration exit rule, a valid-label rate of
19.28\%.

\begin{center}
\begin{tabular}{rrrr}
\toprule
Strategy & Rows & Valid Labels & Valid \% \\
\midrule
K1 & 109,160 & 34,820 & 31.90\% \\
K2 & 109,160 & 10,948 & 10.03\% \\
T1 & 109,160 & 12,510 & 11.46\% \\
T2 & 109,160 & 38,312 & 35.10\% \\
full4 & 109,160 & 8,641 & 7.92\% \\
\bottomrule
\end{tabular}
\end{center}

The main invalid-label reasons are 326,573 rows with invalid/missing target-date
quotes and 113,991 rows with no pre-expiration quote day under the current
horizon rule. This is acceptable for a first training dataset, but the final
paper backtest should not silently drop these rows after trade selection.

\subsection*{Step 4: Leakage-Safe Baseline Backtest Accounting}

The backtest code was patched so trade selection no longer prefilters the test
universe on \texttt{is\_valid\_label}. The model now scores the full
point-in-time candidate universe first; selected trades with missing or invalid
future exits are counted separately from executable trades. Label construction
was also patched to retain entry gross exposure for invalid-exit rows, since
entry exposure is known at decision time.

A first LightGBM baseline was trained only to test the new path, not as a final
research result. The train/validation split used the existing project dates:
2019--2022 for training and 2023 for validation. The model used 42,266 training
rows and 4,843 validation rows, with validation MAE 0.3598 and RMSE 0.6309.

On the 2024 test split, the leakage-safe baseline selected 54 trades. Only 11
had executable realized exits under the current strict target-date quote rule;
43 selected rows had invalid or unavailable exits. Executable-trade P\&L was
\$-103.64 at the current normalized gross-exposure convention, with average
return -0.94\%, median return -0.12\%, and win rate 45.45\%. This result should
be interpreted as a pipeline diagnostic. It shows that the next publishable
step is to define a stronger point-in-time execution and exit policy before
claiming profitability.

\subsection*{Step 5: Three-Month Rolling ML Backtests}

The one-shot train/validation/test split was replaced with an expanding-window
rolling experiment. Each test period is a 3-month calendar window. For each
window, the model trains only on valid realized-return labels from dates before
the test window, then scores all candidate strategy rows in the next 3-month
window. Early windows are skipped until at least 500 prior valid labels exist.
This is not data removal; it is the minimum needed for honest out-of-sample ML
training.

Two LightGBM rolling variants were run:

\begin{enumerate}
  \item \textbf{Strict gates}: prediction $>$ 0, normalized severity $\geq$
  0.0005, max leg spread $\leq$ 25\%, average leg spread $\leq$ 15\%, positive
  volume, positive open interest.
  \item \textbf{All-cleaned gates}: prediction $>$ 0, normalized severity
  $\geq$ 0.0001, max and average spread allowed up to the cleaned-data 50\%
  limit, no volume/OI minimum. This scores and admits much more of the cleaned
  candidate universe, while still using only rows that passed the original
  quote-quality cleaning.
\end{enumerate}

\begin{center}
\begin{tabular}{lrrrrr}
\toprule
Variant & Completed Windows & Test Rows & Selected & Executable & Net P\&L \\
\midrule
Strict gates & 36 & 543,080 & 1,350 & 400 & \$120,523 \\
All-cleaned gates & 36 & 543,080 & 10,279 & 1,005 & \$903,499 \\
\bottomrule
\end{tabular}
\end{center}

The strict run selected only 2.49 trades per 1,000 scored test rows. The
all-cleaned run selected 18.93 trades per 1,000 scored test rows. This explains
the earlier low trade count: the issue was conservative trade gating, not that
the dataset or rolling pipeline ignored most rows.

However, the all-cleaned run also exposes the next major research caveat:
only 1,005 of 10,279 selected trades currently have executable realized exits
under the strict target-date quote rule. The remaining selected rows are mostly
\texttt{invalid\_quote} exit cases. Therefore, the promising P\&L numbers should
be treated as signal diagnostics, not final profitability claims, until the
exit policy is made more realistic and point-in-time.

\begin{center}
\begin{tabular}{rrrrr}
\toprule
Year & Test Rows & Selected & Executable & Net P\&L \\
\midrule
2020 & 133,265 & 1,010 & 221 & \$145,092 \\
2021 & 17,630 & 874 & 69 & \$69,284 \\
2022 & 54,525 & 1,663 & 95 & \$15,818 \\
2023 & 25,300 & 1,049 & 135 & \$154,659 \\
2024 & 68,215 & 1,463 & 164 & \$162,835 \\
2025 & 189,550 & 2,615 & 165 & \$208,510 \\
2026 & 45,270 & 1,008 & 60 & \$143,336 \\
\bottomrule
\end{tabular}
\end{center}

The rolling outputs are saved under:

\begin{itemize}
  \item \texttt{reports/rolling\_3m\_lightgbm/}
  \item \texttt{reports/rolling\_3m\_lightgbm\_all\_cleaned\_gates/}
\end{itemize}

\subsection*{Step 6: Exact 3/1/2 Walk-Forward Protocol}

The previous rolling test used expanding historical training. A stricter
walk-forward protocol was then added to match the requested design:

\[
\text{Train months }1,2,3
\quad\rightarrow\quad
\text{Validate month }4
\quad\rightarrow\quad
\text{Test months }5,6.
\]

The window advances by 2 months so the test windows do not overlap. The
validation month is used to choose the prediction threshold before testing.
The run starts from June 2013, the first month with detected labels. Early
windows are still included in the audit table, but are skipped when a 3-month
training slice has fewer than 20 valid labels or the validation month has no
valid labels. This is an honest ML limitation caused by sparse early violations,
not a data-loading failure.

Important implementation facts:

\begin{itemize}
  \item The label horizon is currently 5 trading days, with fallback to the last
  quote before the earliest selected expiration.
  \item Transaction cost is included as \$0.65 per contract per side; each leg
  pays a round-trip fee of $2 \times 0.65 \times |\text{quantity}|$.
  \item The model uses 104 numeric features and 5 categorical features. It uses
  severity, forward-adjusted strike information, leg-level quotes, IV/Greeks,
  liquidity, moneyness, IV slope/term features, delta gaps, and persistence
  features.
  \item Exact option symbols are \emph{not} used as ML features. They are used
  only as stable quote-lookup keys for exits.
\end{itemize}

The exact 3/1/2 all-cleaned-gates run produced:

\begin{center}
\begin{tabular}{rrrrr}
\toprule
Completed Windows & Test Rows & Selected & Executable & Net P\&L \\
\midrule
55 & 418,155 & 6,027 & 956 & \$268,393 \\
\bottomrule
\end{tabular}
\end{center}

\begin{center}
\begin{tabular}{rrrrr}
\toprule
Year & Test Rows & Selected & Executable & Net P\&L \\
\midrule
2020 & 12,725 & 557 & 108 & \$58,636 \\
2021 & 34,860 & 461 & 40 & \$37,964 \\
2022 & 37,120 & 918 & 112 & \$3,279 \\
2023 & 25,575 & 693 & 97 & \$38,282 \\
2024 & 73,530 & 1,157 & 172 & \$66,941 \\
2025 & 184,800 & 1,668 & 348 & \$78,528 \\
2026 & 43,950 & 441 & 48 & \$-16,446 \\
\bottomrule
\end{tabular}
\end{center}

The exact 3/1/2 protocol is more faithful to short rolling model adaptation,
but less stable than expanding-window training because only 3 months of labels
are used for each model fit. The expanding all-cleaned run produced higher
aggregate P\&L, while the exact 3/1/2 run is more conservative and structurally
closer to the requested research design.

\begin{center}
\begin{tabular}{lrrrr}
\toprule
Protocol & Test Rows & Selected & Executable & Net P\&L \\
\midrule
Expanding strict & 543,080 & 1,350 & 400 & \$120,523 \\
Expanding all-cleaned & 543,080 & 10,279 & 1,005 & \$903,499 \\
Exact 3/1/2 all-cleaned & 418,155 & 6,027 & 956 & \$268,393 \\
\bottomrule
\end{tabular}
\end{center}

\begin{center}
\includegraphics[width=0.92\textwidth]{walkforward_3_1_2_lightgbm_all_cleaned_gates/walkforward_yearly_pnl_pil.png}
\end{center}

\begin{center}
\includegraphics[width=0.92\textwidth]{walkforward_3_1_2_lightgbm_all_cleaned_gates/walkforward_selected_vs_executable_pil.png}
\end{center}

The exact walk-forward outputs are saved under:

\begin{itemize}
  \item \texttt{reports/walkforward\_3\_1\_2\_lightgbm\_all\_cleaned\_gates/}
  \item \texttt{walkforward\_3\_1\_2\_summary.csv}
  \item \texttt{walkforward\_3\_1\_2\_yearly\_summary.csv}
  \item \texttt{walkforward\_3\_1\_2\_window\_table.csv}
  \item \texttt{walkforward\_3\_1\_2\_cumulative\_pnl.csv}
  \item \texttt{protocol\_comparison.csv}
\end{itemize}

\section*{June 9 Moneyness Comparison Extension}

The pipeline was extended from the original OTM-only call universe into a
controlled moneyness comparison. This was done to test whether the TP2-inspired
signal is truly strongest in OTM calls, or whether ATM/ITM/all-call universes
provide stronger publishable evidence. The raw Archive-2 files were not
modified. A new derived splitter was added:
\texttt{scripts/build\_archive2\_moneyness\_universes.py}.

The moneyness buckets are defined using strike divided by raw SPY close:

\begin{itemize}
  \item ITM calls: $K/S < 0.98$.
  \item ATM calls: $0.98 \leq K/S \leq 1.02$.
  \item OTM calls: $K/S > 1.02$.
  \item All calls: all cleaned SPY calls, used as a combined benchmark.
\end{itemize}

Each bucket was processed independently through the same forward-adjusted TP2
detector, the same five strategy-label construction, and the same exact 3/1/2
walk-forward LightGBM protocol. The exit quote panel remained the broad
all-cleaned call panel, so an exit was not lost merely because a contract moved
between moneyness regions after entry.

\begin{center}
\begin{tabular}{lrrrrrrr}
\toprule
Universe & Clean Rows & Violations & Valid Labels & Windows & Selected & Executable & Net P\&L \\
\midrule
ITM & 90,760 & 1,199 & 45.9\% & 2 & 19 & 2 & \$42 \\
ATM & 423,116 & 5,555 & 40.7\% & 38 & 774 & 365 & \$7,255 \\
OTM & 290,933 & 92,458 & 18.4\% & 51 & 4,780 & 780 & \$213,059 \\
All calls & 804,809 & 115,803 & 21.5\% & 55 & 6,930 & 1,184 & \$362,829 \\
\bottomrule
\end{tabular}
\end{center}

\begin{center}
\includegraphics[width=0.92\textwidth]{figures/moneyness_walkforward_comparison.png}
\end{center}

The comparison supports a stronger paper framing than an OTM-only study. ITM
calls are too sparse as a standalone ML universe under the strict $K/S < 0.98$
definition. ATM calls are executable and positive, but materially weaker than
OTM. OTM remains the cleanest theoretically motivated universe and produces a
large positive result. The all-call benchmark produces the largest aggregate
P\&L, but it also has the largest number of invalid selected trades and mixes
OTM, ATM, and ITM mechanics. Therefore, the preferred research claim is not
``all calls are better'' or ``only OTM matters.'' The stronger claim is that
TP2-inspired trading performance is moneyness-dependent, with OTM as the
primary clean universe and all-calls as a high-capacity benchmark requiring
stricter execution robustness.

The new comparison artifacts are saved at:

\begin{itemize}
  \item \texttt{data/processed/archive2\_moneyness/}
  \item \texttt{data/features/archive2\_moneyness/}
  \item \texttt{data/labels/archive2\_moneyness\_strategy\_training\_rows\_*\_2013\_2026.parquet}
  \item \texttt{reports/moneyness\_walkforward\_3\_1\_2/}
  \item \texttt{reports/tables/moneyness\_pipeline\_walkforward\_summary.csv}
  \item \texttt{reports/tables/moneyness\_walkforward\_yearly\_summary.csv}
\end{itemize}

\section*{June 9 Mixed-Rectangle Moneyness Check}

One important issue with pre-splitting the universe by moneyness is that a TP2
rectangle has four option legs. A detected rectangle can therefore mix ATM and
OTM legs, or ITM and ATM legs. The independent ITM/ATM/OTM runs above are valid
``within-bucket'' tests, but they intentionally exclude mixed-moneyness
rectangles.

To check whether the ML pipeline was missing useful mixed rectangles, the
detector was run on the all-call universe first. Each detected rectangle was
then labelled by the entry-time moneyness of its four legs using
\texttt{scripts/add\_archive2\_rectangle\_moneyness.py}. These moneyness
labels are point-in-time features only; they do not use future exits or future
profit information.

\begin{center}
\begin{tabular}{lrr}
\toprule
Rectangle Type & Violations & Share of All-Call Violations \\
\midrule
Pure OTM & 89,014 & 76.9\% \\
Mixed ATM/OTM & 16,660 & 14.4\% \\
Pure ATM & 5,580 & 4.8\% \\
Mixed ITM/ATM & 1,972 & 1.7\% \\
Pure ITM & 1,226 & 1.1\% \\
Mixed ITM/OTM & 1,076 & 0.9\% \\
Mixed ITM/ATM/OTM & 275 & 0.2\% \\
\midrule
Total & 115,803 & 100.0\% \\
\bottomrule
\end{tabular}
\end{center}

The category-aware all-call label file contains 579,015 strategy rows, exactly
five candidate strategies for each of the 115,803 violations. The ML feature
audit confirmed that target and future columns such as
\texttt{realized\_return}, \texttt{net\_pnl}, \texttt{exit\_date},
\texttt{is\_valid\_label}, and contract identifiers are excluded from the model
feature set. The newly added rectangle moneyness fields are retained because
they are known at entry.

\begin{center}
\begin{tabular}{lrrrr}
\toprule
Rectangle Type & Selected & Executable & Invalid Selected & Net P\&L \\
\midrule
Pure OTM & 4,968 & 694 & 4,274 & \$380,026 \\
Mixed ATM/OTM & 1,342 & 213 & 1,129 & \$17,344 \\
Pure ATM & 661 & 287 & 374 & -\$1,914 \\
Pure ITM & 54 & 9 & 45 & -\$491 \\
Mixed ITM/ATM & 57 & 20 & 37 & -\$771 \\
Mixed ITM/OTM & 8 & 1 & 7 & -\$177 \\
Mixed ITM/ATM/OTM & 1 & 1 & 0 & \$195 \\
\midrule
Total & 7,091 & 1,225 & 5,866 & \$394,212 \\
\bottomrule
\end{tabular}
\end{center}

\begin{center}
\includegraphics[width=0.92\textwidth]{figures/rectangle_moneyness_category_pnl.png}
\end{center}

The category-aware all-call run is stronger than the earlier all-call run that
did not explicitly expose rectangle moneyness to the model:

\begin{center}
\begin{tabular}{lrrrr}
\toprule
Run & Selected & Executable & Invalid Selected & Net P\&L \\
\midrule
All calls without explicit rectangle category & 6,930 & 1,184 & 5,746 & \$362,829 \\
All calls with rectangle category features & 7,091 & 1,225 & 5,866 & \$394,212 \\
\bottomrule
\end{tabular}
\end{center}

This resolves the main concern about splitting before training. Pure OTM
rectangles remain the dominant economic source of profit, while mixed ATM/OTM
rectangles add a meaningful secondary contribution. Pure ATM is more executable
but negative in this run, and ITM-related rectangles are too small or weak to
support a standalone claim. For the paper, the clean framing is: detect on all
calls, label the rectangle moneyness structure, and report that TP2-inspired
profitability is concentrated in OTM-heavy rectangles rather than uniformly
across all call moneyness regions.

The remaining execution caveat is the high invalid-selected count. The current
walk-forward code scores the full test set and reports invalid selected trades
instead of filtering them away before selection, so there is no future-label
selection leakage in the test accounting. However, before paper submission the
strategy should still add stricter liquidity and exit-policy robustness checks,
including flexible exits rather than only the current five-trading-day target
exit.

\section*{June 9 Flexible Exit Robustness Check}

The fixed five-trading-day exit is useful as a clean baseline, but it is not a
complete trading policy. If a selected contract cannot be quoted cleanly on the
target exit date, the backtest cannot assume an automatic fifth-day sale. A
stronger paper needs to report this explicitly and test a more realistic exit
rule.

A separate flexible-exit label builder was therefore added:
\texttt{src/tp2rr2/flexible\_labels.py} with runner
\texttt{scripts/run\_flexible\_exit\_labels.py}. The original fixed-exit label
files were not overwritten. The flexible policy is:

\begin{enumerate}
  \item Enter only after a TP2 violation is observed.
  \item For each future trading day, check same-contract quotes for the four
  rectangle legs.
  \item Exit on the first valid day when the rectangle no longer violates by the
  entry threshold, i.e. normalized severity falls to $\leq 0.0001$.
  \item If no valid reversion appears before the relevant expiration, exit on
  the latest valid pre-expiration quote for the selected strategy legs.
  \item If no valid exit quote exists, mark the selected row as invalid rather
  than silently dropping it before selection.
\end{enumerate}

This raised the label validity rate from roughly 21.5\% under the fixed
five-day all-call labels to 42.0\% under flexible exits. The flexible label file
contains the same 579,015 strategy rows and the same 115,803 violation IDs as
the fixed-exit file.

During this check, two future-known label fields were also removed from the ML
feature set: \texttt{actual\_holding\_days} and
\texttt{exit\_normalized\_severity}. The corrected model exclusion list now
drops target, exit, validity, and identifier fields before training. The
fixed-exit category-aware model was rerun after this correction, so the
comparison below is leakage-clean.

\begin{center}
\begin{tabular}{lrrrrr}
\toprule
Exit Policy & Completed Windows & Selected & Executable & Invalid Selected & Net P\&L \\
\midrule
Fixed 5-day, corrected no-leakage & 55 & 7,304 & 1,192 & 6,112 & \$384,129 \\
Flexible reversion/fallback & 57 & 5,597 & 2,309 & 3,288 & \$195,221 \\
\bottomrule
\end{tabular}
\end{center}

\begin{center}
\includegraphics[width=0.92\textwidth]{figures/exit_policy_walkforward_comparison.png}
\end{center}

The flexible exit materially improves executability: executable selected trades
rise from 16.3\% to 41.3\% of selected trades, and invalid selected trades fall
from 83.7\% to 58.7\%. However, aggregate P\&L falls from about \$384k to
\$195k. The selected flexible-exit trades split as follows: 1,232 reversion
exits produced about \$276.8k, while 1,077 latest-valid pre-expiration fallback
exits lost about \$81.6k. This means the reversion idea is economically useful,
but the fallback component needs stricter gates before it can be treated as a
production-quality trading rule.

The paper-safe interpretation is therefore not that flexible exits dominate the
fixed baseline. The stronger claim is that the TP2 signal survives a
leakage-clean ML rerun, and that a realistic reversion/fallback exit policy
substantially improves execution validity while revealing where additional
liquidity and fallback-exit controls are needed. For the main paper tables, the
corrected fixed-exit result should be reported as the high-P\&L baseline and
the flexible-exit result should be reported as an execution robustness
specification.

\section*{Active Pipeline Versus Legacy Artifacts}

The active Archive-2 research pipeline is now:

\begin{enumerate}
  \item \texttt{scripts/run\_ingest\_archive2.py}
  \item \texttt{src/tp2rr2/archive2\_ingest.py}
  \item \texttt{scripts/build\_archive2\_moneyness\_universes.py}
  \item \texttt{scripts/add\_archive2\_rectangle\_moneyness.py}
  \item \texttt{scripts/run\_detect\_archive2\_forward.py}
  \item \texttt{src/tp2rr2/detect\_forward\_tp2.py}
  \item \texttt{scripts/run\_labels.py}
  \item \texttt{src/tp2rr2/flexible\_labels.py}
  \item \texttt{scripts/run\_flexible\_exit\_labels.py}
  \item \texttt{scripts/run\_walkforward\_3\_1\_2.py}
\end{enumerate}

Older phase reports, SHAP reports, and earlier JSON-archive outputs should be
treated as legacy artifacts from the previous project stage. They should not be
used for the current ICAIF Archive-2 result unless explicitly re-run on the new
forward-adjusted Archive-2 pipeline. The literature PDFs should remain in place.

Remaining implementation caveats before paper-quality ML/backtesting:

\begin{itemize}
  \item The ingestion should make the \texttt{quote\_time} timezone convention
  explicit, because the raw README describes UTC instants while the retained
  timestamps appear as 15:55--15:59 after local reading.
  \item The flexible fallback exit should be tightened. Reversion exits are
  profitable in the current robustness run, but latest-valid pre-expiration
  fallback exits are negative and should require additional liquidity,
  moneyness, and holding-period controls.
\end{itemize}

\section*{Next Updates}

Future project-affecting changes should be documented side by side in report
files. The next report updates should cover:

\begin{itemize}
  \item Create a broad call-only exit quote panel for label construction.
  \item Rebuild 2013--2026 strategy labels using OTM-only detection features
  and the broad exit quote panel.
  \item Update backtesting so trade selection is point-in-time and no longer
  prefilters on future label availability.
  \item Define and test a stronger execution policy for selected trades whose
  target-date exit quotes are missing, including DTE/liquidity gates and a
  transparent fallback exit convention.
  \item Add rolling-window result tables and compare strict gates against
  all-cleaned gates in the paper draft.
  \item Add dividend and ex-dividend risk filters for American-call assignment
  and early-exercise risk.
  \item Run robustness checks for strike caps, severity thresholds, quote-time
  strictness, spread thresholds, open-interest/volume gates, and two-strike
  approximation comparison.
  \item Only after those checks, train machine-learning/deep-learning trade
  selectors and report validation/test results.
\end{itemize}

\end{document}
